\lecture{8}{4. November 2025}{Fri svingning/egensvingning af MDOF-systemer (del III)}

\begin{problem}{8.1}
  Vi returnerer til systemet behandlet i Opgave 7.1. Bevis, at dæmpningsforholdene er givet som
  \[ 
  \xi_1 = 0, \quad \xi_2 = \frac{c}{m \omega_2}
  .\]
\end{problem}

\begin{derivation}
  Fra forelæsningen har vi at $\tilde{C}$ er givet som
  \[ 
  \tilde{C}_{jj} = 2 \xi_{j} \tilde{m}_{jj} \omega_j
  .\]
  Heri kan dæmpningsforholdet $\xi_{j}$ isoleres som
  \[ 
  \xi_{j} = \frac{\tilde{C}_{jj}}{2 \tilde{M}_{jj} \omega_j}
  .\]
  Hvori $\tilde{C}$ er fundet i opgave 7.1 til
  \[ 
  \tilde{C} = \begin{bmatrix}
  0 & 0\\
  0 & 4c\\
  \end{bmatrix}
  .\]
  
  Fra Opgave 6.1 har vi hhv:
  \[ 
  \omega_1 = \sqrt{\frac{k}{m}}, \quad \omega_2 = \sqrt{5}\omega_1
  .\]
  De modale masser bestemmes som
  \begin{align*}
    \tilde{M} &= \Phi^{\intercal} M \Phi \\
              &= \begin{bmatrix}
              1 & 1\\
              1 & -1\\
              \end{bmatrix} \begin{bmatrix}
              m & 0\\
              0 & m\\
              \end{bmatrix} \begin{bmatrix}
              1 & 1\\
              1 & -1\\
              \end{bmatrix} \\
              &= 2m \begin{bmatrix}
              1 & 0\\
              0 & 1\\
              \end{bmatrix}
  .\end{align*} 
  Dermed har vi at $\tilde{M}_{11} = \tilde{M}_{22} = 2m$, hvorfor vi får:
\end{derivation}
\finalresult{
  \begin{aligned}
    \xi_1 &= \frac{0}{2 \cdot 2m \cdot \omega_1} = 0 \\
    \xi_2 &= \frac{4c}{2 \cdot 2m \cdot \omega_2} = \frac{c}{m \omega_2}
  .\end{aligned}
}

\begin{problemwithimage}{p8_2}{8.2}
  Et eksperiment har--under antagelsen om klassisk dæmpning--resulteret i de nedenfor viste egenværdier (og deres komplekst konjugerede), der ligger på en ret linje fra origo.
  Det oplyses, at $\lambda_1 = \num{-0,900} + \num{8,956}i $ og $\lambda_2 = \num{-2,627} + \num{26,142} i$.
  (a) Bestem de fem dæmpningsforhold.

  (b) Bestem de to første udæmpede vinkelegenfrekvenser.

  Det oplyses nu, at
  \[ 
  K = \num{1000} \begin{bmatrix}
  2 & -1 & 0 & 0 & 0\\
  -1 & 2 & -1 & 0 & 0\\
  0 & -1 & 2 & -1 & 0\\
  0 & 0 & -1 & 2 & -1\\
  0 & 0 & 0 & -1 & 1\\
  \end{bmatrix}, \qquad M = I
  .\]
  
  (c) Bestem dæmpningsmatricen, $C \in \mathbb{R}^{5 \times 5}$, ved brug af Rayleighs dæmpningsmodel baseret på de to første egensvingningskonfigurationer.

  (d) Hvor store fejl yder Rayleighs dæmpningsmodel på $\xi_3, \xi_4$ og $\xi_5$?
\end{problemwithimage}

\begin{derivation}
  Fra Forelæsning 7 har vi at der for klassisk dæmpning gælder at:
  \[ 
  \left\{ \lambda_j, \lambda_{j+n} \right\} = - \xi_j \omega_j \pm i \omega_j \sqrt{1 - \xi_j^2}, \text{ hvor } \xi_j \equiv \frac{\tilde{C}_{jj}}{2 \tilde{M}_{jj} \omega_j} \text{ for } \omega_j > 0
  .\]
  Heraf ses at hældningen på den rette linje der går igennem alle egenværdierne er
  \begin{equation} \label{eq:hældning}
    m = \frac{\Im(\lambda_j)}{\Re(\lambda_j)} = \frac{\omega_j \sqrt{1 - \xi_j^2}}{- \xi_j \omega_j} = - \frac{\sqrt{1- \xi^2_j}}{\xi_j}
  \end{equation}
  Her ses at hældningen udelukkende afhænger af dæmpningsforholdet.
  Idet hældningen er konstant må dæmpningsforholdene ligeledes være konstante
  \[ 
  \xi_1 = \xi_2 = \xi_3 = \xi_4 = \xi_5
  .\]
  I \Cref{eq:hældning} kan dæmpningsforholdet isoleres som:
  \begin{align*}
    m \xi_j &= \sqrt{1 - \xi_j^2} \\
    m^2 \xi_j^2 &= 1 - \xi_j^2 \\
    m^2 \xi_j^2 + \xi_j^2 &= 1 \\
    \xi_j^2 \left( m^2 + 1 \right) &= 1 \\
    \xi_j &= \sqrt{\frac{1}{m^2 + 1}} \\
       &= \sqrt{\frac{1}{\left( \frac{\Im(\lambda_j)}{\Re(\lambda_j)} \right)^2 + 1}} \\
    \xi_1 &= \sqrt{\frac{1}{\left( \frac{\num{8,96}}{\num{-0,900}} \right)^2 +1}} \\
          &= \num{0,1}
  .\end{align*}

  Vinkelegenfrekvenserne kan findes ud fra $\lambda$-værdierne som:
  \begin{align*}
    \lambda_j &= - \xi_j \omega_j \pm i \omega_j \sqrt{1 - \xi^2_j} \\
    |\lambda| &= \sqrt{\left( -\xi_j \omega_j \right)^2 \pm \left( i \omega_j \sqrt{1 - \xi^2_j} \right)^2} \\
        &= \sqrt{\xi_j^2 \omega_j^2 + \omega_j^2 - \xi_j^2 \omega_j^2} \\
        &= \omega_j
  .\end{align*}
  Altså kan den $j$'te vinkelegenfrekvens findes som længden af den $j$'te komplekse egenværdi som
  \begin{align*}
    \omega_1 &= \sqrt{\num{0,9}^2 + \num{8,96}^2} = \num{9,00508754} \\
    \omega_2 &= \sqrt{\num{2,63}^2 + \num{26,1}^2} = \num{26,2321729942} 
  .\end{align*}

  Rayleighs dæmpningsmodel foreskriver at:
  \[ 
  C = \alpha M + \beta K
  \]
  hvor $M$ og $K$ i dette tilfælde er givet.
  Vi kan skrive:
  \begin{align*}
    \begin{bmatrix}
    \xi_1\\
    \xi_2\\
    \end{bmatrix} &= \begin{bmatrix}
    \frac{1}{2\omega_1} & \frac{\omega_1}{2}\\
    \frac{1}{2\omega_2} & \frac{\omega_2}{2}\\
    \end{bmatrix} \begin{bmatrix}
    \alpha\\
    \beta\\
    \end{bmatrix} \\
    \begin{bmatrix}
    \alpha\\
    \beta\\
    \end{bmatrix} &= \begin{bmatrix}
    \frac{1}{2\omega_1} & \frac{\omega_1}{2}\\
    \frac{1}{2\omega_2} & \frac{\omega_2}{2}\\
    \end{bmatrix}^{-1} \begin{bmatrix}
    \xi_1\\
    \xi_2\\
    \end{bmatrix} \\
    &= \begin{bmatrix}
    \frac{1}{2 \cdot \num{9,00508754}} & \frac{\num{9,00508754}}{2}\\
    \frac{1}{2 \cdot \num{26,2321729942}} & \frac{\num{26,2321729942}}{2}\\
    \end{bmatrix}^{-1} \begin{bmatrix}
    \num{0,1} \\
    \num{0,1} \\
    \end{bmatrix} \\
    &= \begin{bmatrix}
    \num{1,341} \\
    \num{0,006} \\
    \end{bmatrix}
  .\end{align*}
  Dermed får vi altså:
  \[ 
  C = \num{1,341} M + \num{0,006} K \approx \begin{bmatrix}
  \num{13,341}  & -6 & 0 & 0 & 0\\
  -6 & \num{13,341}  & -6 & 0 & 0\\
  0 & -6 & \num{13,341}  & -6 & 0\\
  0 & 0 & -6 & \num{13,341}  & -6\\
  0 & 0 & 0 & -6 & \num{7,341} \\
  \end{bmatrix}
  .\]  
\end{derivation}
